\section{Data and task definition}
\paragraph{Sources.} Incident records come from the Toronto Police Service open-data Major Crime Indicator (MCI) files for robbery, auto theft, break and enter, and theft over. Neighbourhood features come from a City of Toronto census neighbourhood profile joined by neighbourhood: population, household size, rental properties, median and average individual income, three ratios of workers to population or labour force, and the difference between median and average individual income. These features are fixed within a neighbourhood and over time.

\paragraph{Definition.} The unit is one incident record. The positive class is robbery; the negative class is auto theft, break and enter or theft over. Assault is a Major Crime Indicator but is not a property crime and is not in the inputs; bicycle theft, homicide and shootings are excluded by rule. The window is occurrence years 2014 to 2023, since reporting in these files starts in 2014. Rows need a valid date, an hour of day, a premises type, coordinates inside the Toronto bounding box and a neighbourhood other than ``not a neighbourhood''; exact duplicates within a category are removed. Events labelled both robbery and another type are removed (\resG{log.cross_category_conflict_events_dropped} events), because their class is ambiguous. Table~\ref{tab:flow} gives the sample flow.

\begin{table}[t]
\centering\small
\caption{Sample flow and class counts. Each removal is a row count; the steps add up to the unified table.}\label{tab:flow}
% Generated by tools/gen_paper.py from results/. Do not edit by hand.
\begin{tabular}{@{}lr@{}}
\toprule
Step & Rows \\
\midrule
Rows in the four Toronto Police files & 174,993 \\
Removed: occurrence outside 2014 to 2023, or invalid date & 332 \\
Rows in the unified table & 159,072 \\
Removed: no neighbourhood (NSA) & 16 \\
Rows used by the models & 159,056 \\
\midrule
\multicolumn{2}{@{}l}{\textit{Model rows by class}} \\
Robbery & 25,657 \\
Auto Theft & 52,223 \\
Break and Enter & 69,108 \\
Theft Over & 12,068 \\
\bottomrule
\end{tabular}

\end{table}

\paragraph{Composition shifts over time.} The class mix is not stable (Figure~\ref{fig:comp}). Robbery incidents fell from \resG{category_by_year.2014.Robbery} in 2014 to \resG{category_by_year.2023.Robbery} in 2023, while auto thefts rose from \resG{category_by_year.2020.Auto Theft} in 2020 to \resG{category_by_year.2023.Auto Theft} in 2023. Robbery's share of the four-crime total therefore fell from \resG{prevalence_by_year.2014.share} to \resG{prevalence_by_year.2023.share}, partly because the denominator grew. Any model of this outcome inherits that drift.

\begin{figure}[t]
\centering
\includegraphics[width=\linewidth]{build/figures/fig1.pdf}
\caption{Incidents per year by crime type (left) and the share of incidents that are robberies (right), 2014--2023. Auto theft rises sharply after 2021, and robbery is lower in 2023 than in 2014, so the robbery share falls partly through the denominator.}\label{fig:comp}
\end{figure}
